% ============================================================
%  第1章 异构计算简介
%  源文件：content/chapter1/
% ============================================================


% ---- 1.1-chinese.md ----
\section{异构计算介绍}

异构计算包括串行处理和并行处理。在异构计算环境下，构成应用程序的任务会被映射到系统中最合适的处理设备上。这样，系统中的程序就能并行或并发的执行，从而提高性能和节省功耗。开放计算语言(OpenCL)是一种适合在支持异构开发环境中进行编程的语言。为了帮助读者理解OpenCL2.0中的新特性，我们将从介绍异构和并行计算开始。之后，我们会给OpenCL以合理的定位，基于该定位来讨论OpenCL中的异构编程模型。

当前，异构计算的环境越来越多元化，为了尽可能发挥多核微处理器、中央处理单元群(CPUs)、数字型号处理器(DSP)、可擦写配置的硬件(FPGA，可编程门阵列)，还有图形处理单元(GPUs)。在目前这些异构环境下，需要有能够兼容处理这些异构环境的架构，这样的架构也需要将软件任务映射到异构设备上。不过，这种的架构会给当前的编程体系带来许多挑战。

异构应用通常会对负载行为进行混合，从控制密集型(比如，搜索、排序和分析)到数据密集型(比如，图像处理、模拟和建模，以及数据挖掘)。这些任务也具有计算密集型(比如，迭代函数、数值方法和金融建模)的特点，这种应用的整体吞吐量对底层硬件的计算能力很大。每种负载类型通常都会以最高效的执行方式，执行在指定的硬件架构上。不存在一种设备能够很好的执行所有类型的负载。例如，控制密集型应用在超标量CPU上运行最为合适，CPU的分支预测机制很适合执行这样的应用。不过，数据密集型应用执行在支持向量处理架构的处理器下就会更快的运行。其能同时对多个数据项进行操作，并且这种操作是并行的。



% ---- 1.2-chinese.md ----
\section{本书目的}

本书第一版是同类书中第一本以适合课堂的方式介绍 OpenCL 编程的书。第二版中，我们更新了OpenCL 1.2标准的相关内容。在本版中，我们考虑到OpenCL 2.0在原有OpenCL 1.x的基础上又很大的改动，并且还考察了更广泛的一类应用。本书旨在教会大家在当前所使用的系统总，使用OpenCL作为编程语言进行并行编程。本书也包含很多例子，这些例子有在CPU和GPU上运行的，当然也有在集成在加速处理单元(APU)上执行的例子。本书的另一个目的，是为编程者们提供一份在当前并行系统下\textbf{合理使用OpenCL设计应用}的指南。并且通过OpenCL编程环境，为开发者实现各种不同的抽象和特性。为让读者更好的了解OpenCL，我们准备了一些例子，通过这种方式既能了解OpenCL，又能与相关的优化技术相结合，从而产生激动人心的应用。本书也会讨论一些能为提供相关帮助的工具，这些工具包括分析和调试工具，这些工具可以帮助读者调试工程，不会让读者在编程时有太多的挫败感。如果本书被某位讲师相中，本书也会提供一些演示文档和编程实例，以支持您的工作。更多的相关信息，请访问 \texthref{http://store.elsevier.com/9780128014141}



% ---- 1.3-chinese.md ----
\section{并行思想}

大多数应用开始都是运行在单核上。不过，在高性能计算时代，当系统提供多个计算资源时，这些资源就会用来为应用加速。标准的方法包括“分而治之”和“分散-聚合”的分解方式，为编程者提供一系列高效利用并行资源的方式。“分而治之”是将一个问题递归的划分为数个子问题，直到可用的计算资源能够解决划分后的子问题。“分散-聚合”会先发送一部分输入数据到每个并行资源中，然后将这些输入数据处理后的结果进行收集，再将这些结果合并到一个结果数据集中。以前，划分后的任务大小与计算资源的能力有很大关系。图1.1展示了普遍应用处理排序和矢量相乘的方式，以往的应用会将任务映射到可并行的计算资源中进行加速计算。

\myGraphic{0.85}{images/chapter1/1-1.png}{图1.1 (a)简单的排序示例：分治法的实现，先把列表切分成较短的列表，分别排序，再把排好序的短列表合并起来。(b)矢量-标量乘法：把乘法分散到各单元计算，再用一系列步骤把结果归约相加。}

现代计算系统发展中，并行化和异构化成为一种趋势，这样传统的编程方式就一定会受到冲击。因为CMOS(互补金属氧化物半导体，Complementary Metal Oxide Semiconductor)的功耗和热极限是额定的，微处理器供应商发现硬件的瓶颈很难让设备获取更高的性能，所以转而决定在同一块芯片上集成多个（有时是专用的）处理器。因为硬件供应商们已经这样做了，那么对于编程者来说，就再也没有“免费的午餐”了(The Free Lunch Is Over)。同时，需要编程者抽离应用中并行的部分。编程者还要根据相关算法去划分任务，并且将这些划分好的任务映射到不同的硬件平台上。

过去的10年，并行计算设备的数量和计算能力都在快速的增长。近期，GPU出现在计算领域；今天，其能够使用很低的功耗，并提供很高的计算能力，让如今计算机的计算能力达到了一个新的水平。在实时三维图形渲染（一个高度数据并行的问题）需求的驱动下，GPU在其他方面的进化也十分迅速，并且成为一种完全可编程的设备，能够对数据和任务并行的任务进行处理。硬件制造商如今将CPU和GPU集成在了一块芯片上面，这将是下一代异构计算的趋势。计算敏感和数据敏感的应用都可以看做为一个任务，任务的代码在内核中。内核可运行在GPU上，GPU使用低功耗和高性能完成相应的任务，而CPU之后仅负责对非kernel代码任务的执行。

很多系统，以及大自然和人造世界中的某些现象，都展示着不同类别的并行和并发：

\begin{itemize}
\item 分子动力学——分子之间是互相影响的。
\item 天气和海洋模型——成千上万的波浪和潮汐。
\item 多媒体系统——视觉和音效，成千上万像素和波形。
\item 汽车生产流水线——数百辆汽车正在装配，每辆处于不同的装配阶段，并有多条相同的装配线。
\end{itemize}

并行计算由Almasi和Gottlieb[1]定义，这种计算模式下，很多计算可以同时进行。原理就是将大问题分割成小问题，这样小问题之间没有关联，就可以同时进行解决。这种程度的并行，取决于问题本身固有的性质(还记得大自然中存在着有很多并行的现象吧)，应用这样的原理就可以来处理手头上的问题。对于算法技巧和软件设计来说，了解哪些部分可以并行也是很重要的。我们先对两个简单例子进行讨论，展示其内在的并行计算：1.两个整型数组相乘；2.文本查找。

第一个例子会将两个整型数组A和B拆分后进行相乘，每个拆分数组中都有N个数组成，然后将乘后的结果存储到C数组中。图1.2展示了拆分的方式。串行的C++代码如下清单1.1所示。

\myGraphic{0.85}{images/chapter1/1-2.png}{图1.2 将数组A与数组B中的元素相乘，结果存入数组C。}

\begin{codebox}[cpp]
// 清单 1.1: 数组相乘
for (int i = 0; i < N; ++i){
  C[i] = A[i] * B[i];
}
\end{codebox}

虽然这段代的计算敏感度并不高，但这段代码确实可以并行。低计算敏感度说明该算法的计算部分操作，要少于内存操作。A和B中的元素都是相互独立的，如果要并行化这段代码，可以为 C 的每个元素生成一个独立的执行实例来完成计算。这段代码是可以数据并行的，A和B数组可以经过同样的处理，然后产生C数组中的数据。

我们再可以来看一个通过划分方式，实现简单任务并行的例子。任务指一项要完成或承担的工作，有时用来替代操作系统的术语“进程”。这里我们讨论的是，使用整个任务的部分输入作为一个子项进行操作。任务并行可以进一步推广到流水线上，或是复杂的并行交互上。图1.3展示了一个任务并行在流水线的例子，这个例子使用快速傅里叶变换在频域中对图像进行滤波。

\myGraphic{0.85}{images/chapter1/1-3.png}{图1.3 快速傅里叶变换(FFT)应用中的任务并行：三张输入图像分别在三个独立任务中被处理。}

让我们来看第二个例子。这里的计算是统计某个字符串在一段文本中出现的次数，然后统计查找词在文本中出现的次数(可见图1.4)。假设查找词有N个，那我们可以将查找词和输入字符串进行匹配(相当于一个任务)。每个例子中查找词都会和输入的文本字符串进行匹配操作。这种方式虽然搜索效率不高，但并行度很高。不过，每次输入的字符串都要和这N个查找词进行匹配。每个并行任务都在执行相同的操作。单个比较任务内部还存在进一步并行：逐字符比较呈现出更细的粒度。这两个例子展示了数据和任务级别的并行。

\myGraphic{0.85}{images/chapter1/1-4.png}{图1.4 任务级并行：多个词可以同时比较；图中还展示了词内逐字符比较时更细粒度的并行性。}

匹配数量是确定的，那我们需要将其出现的次数进行相加。同样的，相加操作也可以并行化。这一步中，会介绍“归约”原理，我们将使用高效的方式，利用并行资源来对各个词进行统计。图1.5所示的归约树，就是我们用来处理相加计算的方式，其时间复杂度为O(logN)。

\myGraphic{0.85}{images/chapter1/1-5.png}{图1.5 图1.4中的字符串比较全部完成后，可以用归约网络把匹配的数目累加起来。}

\myRefRule

\begin{myRefs}

[1] G.S.Almamsi, A Goulieb, Highlu Parraleel Computing, Benjamin-Cummings Publishing Co.. Inc.. RedWood City, Ca, USA, 1989

\end{myRefs}



% ---- 1.4-chinese.md ----
\section{并发和并行编程模型}

尝试使用OpenCL进行应用开发之前，需要来讨论一下并行和并发执行模型，以便于我们选择合适的模型。虽然，后面的模型OpenCL都能支持，不过具体的底层硬件可能会有一些限制，这就需要进行一定的尝试。

当我们在讨论并发编程时，就意味着独立的系统在执行多个独立的任务。不过，并发任务并不要求必须同时（并行）执行。例如一个简单的绘图程序，其收到用户使用鼠标和键盘发送过来的输入信号，或更新当前显示的图片。应用可以并发的接收输入信号，并处理与显示相关的对应信号，从而更新显示图像。这些任务可以用并发方式表达，但并不需要并行执行。实际上，这些任务执行在一个CPU上，它们不能并行的执行。在这个例子中，应用或操作系统需要在不同任务间进行切换，确保这些任务都能执行在同一个内核上。

并行化关心的是两个或多个事件同时执行，以提高整体性能。比如如下的赋值语句：

\begin{codebox}[text]
// 代码清单1.2
Step 1: A = B + C
Step 2: D = E + G
Step 3: R = A + D
\end{codebox}

A和D分别在第一和第二步进行赋值，并且这两步之间没有数据依赖，相互独立。第一步和第二步等号左右的数据均不相同，所以第一步和第二步可以并行的执行。第三步数据依赖于第一步和第二步的结果，所以其只能在第一步和第二步执行完成之后执行。

并行编程必须是并发的，不过并发编程不一定并行。虽然很多并发程序可以并行执行，但是互相有依赖的并发任务就不能并行了。比如：交错执行就符合并发的定义，而不能并行的执行。所以，并行是并发的一个子集，并发程序是所有程序集的一个子集。图1.6展示了这三种程序集间的关系。

\myGraphic{0.66}{images/chapter1/1-6.png}{图1.6 并行程序与并发程序之间的关系：并行程序和并发程序都是所有程序的子集。}

接下来的章节中，会介绍一些知名的方式进行异构编程。这里对并发和并行的介绍，是为第2和第3章介绍OpenCL做好铺垫。



% ---- 1.5-chinese.md ----
\section{线程和共享内存}

一个执行的程序(称为进程)可能有很多子程序组成，并且这些子程序都有独立的控制流。当该进程启动时，这些子程序就并发的开始执行。这些子程序可以称之为线程。进程中的所有线程都会共享一些资源(比如：内存、打开的文件、全局变量)，不过他们也有属于自己的资源(比如：堆栈、\href{https://zh.wikipedia.org/zh-cn/\%E8\%87\%AA\%E5\%8A\%A8\%E5\%8F\%98\%E9\%87\%8F}{自动变量})。线程使用全局共享地址空间分配出的变量进行通讯。通讯时需要有同步机制来保证同一个内存区域的内容，不会被多个线程更新。

具有共享内存的系统中，所有的处理器都能访问到同一地址空间(比如：能看到同一份全局内存)。共享内存模型中一个关键的特性，就是不需要编程者去管理数据的移动。在这样的系统中，线程如何去更新全局变量，底层硬件和编程者要达成共识，并遵守相关的访问协议。“相关协议”的学术名称为“内存一致性模型”，内存一致性模型已经被很多高级语言所支持，比如：Java、C/C++11和OpenCL。相关的高级结构有互斥量和信号量，以及“获取-释放”语义。程序员会显式告知底层硬件需要使用哪种同步方式，从而能让硬件更加高效的执行并发任务。

随着处理器数量的增加，底层硬件也要花费很大的精力来对共享内存进行支持和管理。总线的长度(与延迟和功耗相关)，硬件结构上接口的数量，以及共享传输总线的数量都将占用很多资源。当我们尝试增加处理器的数量时，这些额外的硬件设备数量将会呈指数级别增长。我们的系统将变得更复杂，同时也增加了很多的额外开销。这些是多核或多处理器系统的弊端，并且这也限制着一起工作处理器的数量。内存一致性模型下，能够使用的处理器相对较少，这是因为共享总线和相干性协议将会成为性能瓶颈。系统中多一些松散的共享内存，相对来说会好一些；当有成规模的内核在共享内存系统中时，其会让系统变得复杂，并且内核间的互相通讯也要付出很高的代价。大多数多核CPU平台都支持某一种共享内存，OpenCL也可以在支持共享内存的设备上运行。



% ---- 1.6-chinese.md ----
\section{消息通讯机制}

相较共享内存，另一个选择就是消息通讯模型。这个模型中，如果当前正在执行的任务需要某些共享数据参与内部计算，可以通过显式通讯的方式进行获取。这种方式可以在一台物理机上使用，也可以在多台物理机之间使用。任务通过消息通讯的机制对数据进行更新。当其他进程在具有相关数据时，数据传输通常需要一些附加操作，比如：一个发送操作必须与一个接收操作相对应。

这个模型对于编程者来说是透明的，通常包含一组与硬件无关的、用于收发消息的例程库。编程者需要对进程间的通讯进行显式的管理。各种各样的消息传递库在 20 世纪 80 年代时，就已经出现了。消息传递接口(MPI)库经历了时间的磨砺，在当今环境下依旧是很受欢迎的消息传递中间件[2]。不过，不同的MPI实现版本之间有很多的不同，其会让程序的可移植性降低。

\myRefRule

\begin{myRefs}

[2] W. Group. E. Lusk, A. Skjellum, Using MPI: Portable Parallel Programming with the Meaasge-passing Interface, second ed.. MIT Press, 1999

\end{myRefs}



% ---- 1.7-chinese.md ----
\section{并行性的粒度}

不论是在线程中使用共享内存，还在使用消息传递机制，我们都能调整线程的粒度。研究并行计算时，粒度是计算量与通讯量之比(比如，需要在线程间进行同步的变量)。

并行的粒度由算法的特点所决定，选择合适的并行粒度对于编程者来说很重要，合适的粒度能够在当前设备上获得最优的性能。有时对于粒度选择就像是在分块，然后决定哪些数据给予哪个任务，选择了合适的块大小可以就能在并行硬件上获得性能收益。如何选择并行程序的粒度，可以参考下面列出一些经验。

使用细粒度并行时，需要考虑如下几点：

\begin{itemize}
\item 计算量不要过大，这样每个线程都有足够的工作可做。
\item 尽可能减少在数据同步上的开销，这样每个线程能够独立的完成自己的任务。
\item 负载的划分也很重要，因为有大量的独立任务需要并行执行，设计良好的任务调度器都能灵活控制负载，并保证在多任务运行的同时，让线程上的达到均衡。
\end{itemize}

使用粗粒度并行时，需要考虑如下几点：

\begin{itemize}
\item 计算量肯定要高于细粒度并行时的计算量，因为不会像细粒度并行那样，有很多线程同时执行。
\item 编程者使用粗粒度编程时，需要了解应用的整体结构，让粗粒度中的每个线程作为任务，服务于应用。
\end{itemize}

结合给出的建议，考虑一下自己的应用要选择哪种粒度的实现。其实，最合理的粒度取决于所要使用的算法和运行的硬件平台。大多数情况下，当同步和通讯的开销大于计算，那么将并行粒度加大，更有利于控制同步和通讯所需的过高开销。细粒度并行能减少负载不均和访存延迟在性能上的开销(对于GPU来说更是如此，其为粒度之细可以切换线程时达到零开销，同时也能隐藏访存延迟)。

当一个算法需要对大量数据进行同一组操作时，就可将数据视为向量，执行同一操作时，多个数据作为输入，经过向量操作后输出多个数据。这种执行方式就是利用单指令多数据(SIMD)的方式对数据进行处理，可并行硬件可以利用这种执行方式，并行的对不同的数据进行处理。这种粒度的并行与向量的大小有关，通过SIMD执行单元的多数据处理来获得应用性能的提升。

寻找和尝试最好的计算粒度时，是否考虑过为什么不将统一程序中的所有副本拷贝到其他处理单元和节点上，如果能够这样，是不是就不用理会执行单元的执行顺序，并且能让程序高效的运行在一个有着多处理器的共享系统上呢？SIMD模型与单程序多数据(SPMD)模型很相似，与 SIMD 模型不同，SPMD 模型不把同步限制在指令边界上，并且允许这些副本任务或内核能够并发的执行。

\subsection{数据共享与同步}

开发异构软件还有一个很关键的因素——数据共享的程度。像想一下，两个任务没有任何数据需要共享，当运行环境或操作系统中的运算资源充足，那么这两个任务可以并行执行。如果在系统运行第一个任务时，第二个任务需要使用到第一个任务产出的结果。要是存在这样的数据依赖关系，那么就必选要介绍一下如何在并行或并发任务时，进行数据的同步。至少跨越该同步点的并行执行会变得更加困难。当然，这会让编程的难度加大，让异构编程更具有挑战性。

编写并发软件的时候，数据共享和同步扮演着很重要的角色。数据共享就有如下的用途：

\begin{itemize}
\item 一个任务中的输入，依赖于另外一个任务的输出——例如：生产者-消费者模型、流水线执行模型。
\item 需要将中间结果进行综合(比如：归约，还有图1.5里的例子)。
\end{itemize}

理想状态下，可以尝试将应用能并行的部分剥离出来，确保并行的部分没有数据依赖，不过这只是理想状态而已。栅栏和锁有时会用来做数据的同步。虽然，现在只是举了这个例子，不过在后面的章节中我们将在OpenCL中重温这个问题。OpenCL提供相应的机制，来完成主机端和设备端的数据同步，或者是完成任务间的数据同步。

\subsection{共享虚拟内存}

很多异构系统，会将任务分散在不同的设备上执行，并且使用显式同步和通讯机制，来完成不同任务在不同设备上的数据同步。共享虚拟内存设置在硬件和软件层面，能够让不同的设备看到的是同一块内存。这个机制让编程更加简单，并且能节省显式的通讯开销。OpenCL 2.0标准开始支持共享虚拟内存，减少了在程序执行中高昂的通讯开销。同样，这样的机制就不需要在每个设备上都存储一组数据副本，也就可以节省内存的开销。当然，我们会在后面章节中，在与OpenCL上下文的相关章节中，讨论更多共享虚拟内存的细节。

OpenCL2.0提供了三种共享虚拟内存：

\begin{itemize}
\item 粗粒度共享缓存
\item 细粒度共享缓存
\item 系统级细粒度共享缓存
\end{itemize}

使用共享虚拟内存需要OpenCL实现将系统中主机和设备端的地址链接起来。这就允许在OpenCL内核中使用数据结构的指针(比如：链表)，之前版本的OpenCL标准是不支持内核中使用自定义数据结构指针。粗粒度缓存支持通过API的调用，来更新整个缓存的内容。细粒度缓存和系统级细粒度的粒度为字节级，这就意味着无论是主机端，还是设备端，对数据的更新将会立即同步。细粒度共享内存是按内存一致模型中定义好的内存序，在同步点对数据进行同步。第6和第7章中，我们会对这部分内容进行更深层次的探讨。



% ---- 1.8-chinese.md ----
\section{使用OpenCL进行异构计算}

读者们应该对并行原理和异构计算的背景有了一定的了解，下面我们来看看哪些特性在OpenCL中得到了支持。这里我们也来简单的回顾一下OpenCL的历史。

OpenCL是一个异构编程架构，其管理者是非盈利技术组织Khronos Group[3]。OpenCL是一个应用开发框架，在其框架下开发的应用，能够在不同的硬件供应商的设备上运行。第一版的OpenCL(1.0)标准在2008年正式发布，并出现在苹果Mac OSX雪豹系统当中。AMD 在同一时期宣布支持 OpenCL。IBM也在该年宣布，其XL编译器在Power架构的处理器下支持OpenCL。2010年，Khronos Group发布了OpenCL 1.1标准。2011年，发布了OpenCL 1.2标准。本书的第一版介绍了很多OpenCL 1.2的特性。2013年，Khronos Group发布了OpenCL2.0标准，其包含如下新特性：

\begin{itemize}
\item 嵌套并行化
\item 共享虚拟内存
\item 管道内存对象
\item C11原子操作
\item 增强图像支持
\end{itemize}

OpenCL支持多级别的并行化，能够将应用高效的映射到同构或异构、单独或多个CPU或GPU上，以及其他硬件供应商提供的硬件系统当中。OpenCL 同时定义了设备端语言和针对系统内设备的主机端管理层，主机作为对其他设备的管理者。设备端语言被设计用来映射内存系统和执行模型。主机端语言用较低的开销，为复杂的并发程序搭建任务管道。总的来说，就是提供了一种方法，让开发者能更高效的在OpenCL的环境下进行算法设计及实现。

OpenCL支持基于任务和基于数据的并行。OpenCL的内核类似于SPMD模型，并行单元（称为工作项）执行内核的实例，内核实例创建并入队后，会被映射到支持标量和矢量的硬件上运行。硬件平台供应商都在迅速的对OpenCL进行支持，将OpenCL标准中所提到相关内容在硬件层面上给予实现。有了这些供应商，就代表着OpenCL具有很大的市场价值，从手机端与嵌入式平台(ARM，Imageination，MediaTek，Texas Instrument)，到桌面和高性能计算(AMD, Apple, Intel, NVIDIA和IBM)。更多架构的CPU已经支持OpenCL(包括X86, ARM和Power)，以及以吞吐量和向量处理的著称的GPU，还有能细粒度并行的设备，比如FPGA(Alter, Xilinx)。更重要的是，OpenCL是跨平台、全行业支持的，其有良好的编程模型，值的开发者学习和使用。OpenCL将会在越来越多的地方是用到，并且使用和学习的门槛将越来越低。

OpenCL 2.0开始支持共享虚拟内存，后面的章节中会详细探讨。共享虚拟内存是一项很重要的特性，其对能减轻编程者的负担，特别是在类似APU这样使用统一物理内存的系统上。OpenCL 2.0也支持内存一致模型，提供“获取/释放”语义以缓解编程者在不明确的锁上耗费不必要的精力。支持共享内存通讯、管道，还有其他很多OpenCL 2.0的新特性将在后面的章节中一一详细介绍。

\myRefRule

\begin{myRefs}

[3] Khronos Group, OpenCL, 2014. \texthref{https://www.khronos.org/opencl}

\end{myRefs}



% ---- 1.9-chinese.md ----
\section{本书结构}

\begin{itemize}
\item 第1章，介绍并行算法和软件开发的相关原理。其中包括并发、线程和不同粒度的并行化的原理。
\item 第2章，介绍一些支持OpenCL的硬件处理器，包括CPU、GPU和APU。还有不同的指令，比如SIMD和VLIW。该章也会涉及多核处理器的原理，以及面向吞吐量的系统，当然也有更加高级的异构架构。
\item 第3章，开始介绍OpenCL，包括主机端API和设备端C语言。这章会让读者尝试使用OpenCL进行编程。
\item 第4章，深入了解OpenCL编程实例，包括直方图、图像旋转和卷积，并演示一些OpenCL2.0特性，比如管道内存对象。其中的一个例子使用到了OpenCL C++包装器。
\item 第5章，讨论并发和OpenCL编程模型。这章中，我们会讨论内核、工作项，以及OpenCL执行和内存模型。也会展示如何使用OpenCL对任务进行排队和同步。这会让读者更加了解如何书写OpenCL程序，并与OpenCL内存模型进行互动。
\item 第6章，讨论OpenCL主机端的内存模型，包括资源和内存管理。
\item 第7章，继续讨论OpenCL设备端的内存模型。设备端的内存模型将解释执行单元如何从不同的内存地址上对数据进行访问。这章也会更新OpenCL内存一致性模型，包括内存顺序和范围。
\item 第8章，讨论OpenCL在三种完全不同的异构平台上的表现，这三个异构平台包括：(1)AMD FX-8350 CPU (2)AMD Radeon R9 290X GPU (3)AMD A10-7850K APU。这章也会提到内存方面的优化策略。
\item 第9章，提供一个学习的例子，了解一下图像的聚类和搜索。
\item 第10章，介绍如何使用AMD CodeXL对OpenCL程序进行分析和调试。
\item 第11章，介绍\texttt{C++ AMP}，能够让使用者使用C++来发挥硬件并发的能力。
\item 第12章，介绍WebCL，是基于Web的应用使用OpenCL，当硬件支持OpenCL时，能达到加速网页应用的性能。
\item 第13章，讨论如何使用其他高级语言来使用OpenCL，比如：Java、Python和Haskell。
\end{itemize}


